\section{Evidence and Witness Model}\label{sec:evidence}
Every record carries an evidence class (observation, derived, inferred, proof) and an optional status in the specification. Runs and reproducibility checks are observations; the changed-region record and the witness difference are derived; the sufficiency record is derived and marked uncertain. Records are content-addressed and appended to a hash-chained JSONL log. The log holds an evidence package (configuration and $H$, harness digest, environment, policy, inputs, per-input rows with run ids, regions) so that \texttt{verify} can rebuild changed regions from the stored sources, re-derive each input outcome from the stored runs and re-run the decision function, requiring value, level, reasons and witness to match the recorded verdict. A trace function walks a verdict to its witness or sufficiency record, then to the runs and observations behind it. The run-set commitment inside the sufficiency record is a Merkle root; the log itself is a flat hash chain, not a Merkle tree.

\textbf{Integrity status (limited).} The chain is a plaintext, unkeyed SHA-256 chain with content-addressed ids. It detects accidental corruption and inconsistent edits. It is \emph{not tamper-proof}, and it is tamper-evident only in the limited sense above, not against a deliberate writer: anyone who can write the evidence directory can edit records, recompute every id and chain hash and rewrite \texttt{verdict.json} together with its \texttt{log\_head}, and \texttt{verify} will pass, because it checks verdict logic and chain consistency, not who wrote the facts. \texttt{verify} takes recorded observations, per-run status and sandbox capability states as given; it executes nothing. A trust anchor outside the writer's reach (a pinned \texttt{log\_head} or a signature) does not exist and is Phase 3 work. Experiment E7 exercised these properties (Section~\ref{sec:results}, Table~\ref{tab:e7}). The log is plaintext and contains source code, inputs and captured output, without redaction.
